\newcommand{\ord}{\mathrm{ord}}

\newcommand{\GDValo}{\mathrm{DivVal}^{G, \circ}}
\newcommand{\cano}{\mathrm{K}}

\newcommand{\lct}{\mathrm{lct}}

For a $\mathbb{Q}$-Fano variety $X$, K-stability provides a criterion for the existence of Kähler--Einstein metrics, which is used to be known as the Fano Yau-Tian-Donaldson conjecture.
% 
It is proved in \cite{CDS_KE_Fano} for smooth varieties (Fano manifolds), and in \cite{Berman_Kps_QFano_KE, LXZ_fg_val_compu_thre} for $\mathbb{Q}$-Fano varieties. See \cite{Xu_Book} for a detailed study.

The original definition of K-stability is through (Donaldson's) Futaki invariant on test configurations, but the invention of the so-called stability threshold (in \cite{FO_Kst_Kano,BJ_thre_val_kst}) is crucial in further developments.

Verifying K-stability (or estimating the stability thresholds) for an arbitrary $\mathbb{Q}$-Fano variety was challenging until the Abban--Zhuang method appeared in \cite{AZ_Kstab_fano_flag}. However, several classes of varieties were studied prior to that due to their abundance of symmetries, e.g., toric and spherical varieties.

In this paper, we use the idea of the compatible divisors from \cite{AZ_Kstab_fano_flag} to give an alternative and simpler proof for the formula of equivariant stability threshold for toric and spherical varieties.

In \cite{BJ_thre_val_kst}, the authors show that, for a $\mathbb{Q}$-Fano toric variety $X$ with an algebraic torus $T \subseteq X$, the stability threshold $\delta(X)$ (see \Cref{def:delta}) is equal to the log canonical threshold of a $T$-invariant anticanonical $\mathbb{Q}$-divisor, which we denote by $D^T_X$.
% 
Their proof consists of two parts, first prove the equality $\delta(X) = \delta_T(X)$, where $\delta_T(X)$ denotes the $T$-equivariant stability threshold (see \Cref{def:deltaG}), and then get a formula of the threshold $\delta_T(X)$.
% 
The former part simply follows from Zhuang's equivariant criterion \cite{Zhuang_equi_Kst} (see \Cref{lem:toric:delta}); the latter part is simplified to a short argument using compatible divisors (see \Cref{sec:comp:tor}).

% The $D^T_X$ of a $\mathbb{Q}$-Fano toric variety $X$ is exactly the unique $T$-invariant anticanonical $\mathbb{Q}$-divisor ``asymptotically'' compatible with all $T$-invariant divisorial valuations on $X$.

A direct corollary is a rephrasing of the classical criterion for toric varieties,
\begin{corollary}[c.f., \Cref{cor:tor:Kss}]
	The $\mathbb{Q}$-Fano variety $X$ is K-semistable iff the $\mathbb{Q}$-divisor $D^T_X = D_1 + \dots + D_k$, where $D_1, \dots, D_k$ are the $T$-invariant prime divisors on $X$. That is, the $\mathbb{Q}$-divisor $D^T_X$ is identical to the standard choice of the anticanonical divisor for a toric variety.
\end{corollary}

We then find the analogy of the $\mathbb{Q}$-divisor $D^T_X$ on spherical varieties (see \Cref{sec:prelim:sph}). Spherical varieties are the natural generalization of toric varieties, flag varieties and symmetric varieties.

Given a connected reductive group $G$ and a Borel subgroup $B$ of group $G$, a spherical $G$-variety is a $G$-variety that admits an open $B$-orbit. See, for example, \cite{Timashev_Homo_space_equivar_embed} for more properties of spherical varieties. Our main result is to find the $B$-invariant anticanonical $\mathbb{Q}$-divisor $D^B_X$ that computes the $G$-equivariant stability threshold $\delta_G(X)$ (see \Cref{sec:prelim:kstab:GKst}) for a $\mathbb{Q}$-Fano spherical $G$-variety $X$.

\begin{theorem}[c.f., \Cref{thm:sph}]
	Let $X$ be a $\mathbb{Q}$-Fano spherical $G$-variety, the $G$-equivariant stability threshold of $X$ satisfies
	\begin{equation*}
		\delta_G(X) = \inf_{v \in \GDValo_X} {A_X(v) \over v(D^B_X)} = \min_{i = 1, \dots, k} {A_X(v_i) \over v_i(D^B_X)}
	\end{equation*}
	for some $B$-invariant effective anticanonical $\mathbb{Q}$-divisor $D^B_X$ (which is uniquely determined by the variety $X$), where 
	\begin{itemize}
		\item $\GDValo_X$ is the set of non-zero $G$-invariant divisorial valuations on $X$ (see \Cref{sec:prelim:bir:val});
		\item $A_X(v)$ is the log discrepancy of $v$ (see \Cref{sec:prelim:bir:dis}); and
		\item $\{v_i\}_i$ is a prescribed finite subset of the $G$-invariant divisorial valuations on $X$.
	\end{itemize}
\end{theorem}
\begin{remark*}
	The second equality follows from the piecewise $\mathbb{R}_{\geq 0}$-linearity of log discrepancy $A_X$, see \cite{Golota_Delta_Fano_large_aut} (based on \cite[Proposition 5.2]{Pas_sing_sphe}) for the details.
\end{remark*}

As mentioned earlier, by \cite{Zhuang_equi_Kst}, a $\mathbb{Q}$-Fano variety with a group action is K-semistable iff it is equivariantly K-semistable. Therefore, the $\mathbb{Q}$-divisor $D^B_X$ is the key to study the K-stability of spherical varieties.

There have been already several studies on the K-stability of $\mathbb{Q}$-Fano spherical varieties, via geometric analysis (\cite{Delcroix_Kst_Fano_spher}), test configurations (\cite{LLW_wei_Kst_Fano_sph}), and valuative invariants (\cite{Yin_valu_lin_sph}). Our contribution is simplifying \cite[Theorem 4.1]{Yin_valu_lin_sph}, and our formula has the advantage of being free from an \textit{a priori} choice of a $G$-linearization (see \Cref{def:linz}).

\subsection*{Acknowledgements}
I would like to thank Hamid \textsc{Abban} and Johannes \textsc{Hofscheier} for their invaluable guidance and support throughout the study of K-stability and spherical geometry. I am also grateful to Thibaut \textsc{Delcroix} and \textsc{Yin} Chenxi for helpful discussions during the initial stage of this project. Furthermore, the careful reading and constructive comments from the anonymous referee greatly improved the manuscript.

\subsection*{Convention}
Throughout this paper, we work over the field of complex numbers $\mathbb{C}$. A \emph{variety} means a reduced, separated scheme of finite type over $\mathbb{C}$, and we further require them to be irreducible and normal. In particular, all algebraic groups are assumed to be connected.
% 
A \emph{point} of a variety is a closed point (i.e., $\mathbb{C}$-point).
